\section{Discussion}\label{sec:disc}

\subsection{Answer to the research question}
The mask-aware PI-GRU improves on the published baseline for Building 59 by a clear margin: 7\% lower RMSE, 19\% lower MAE and 38\% lower MAPE on complete data, and 20\% lower RMSE than the re-implemented baseline when 40\% of the readings are lost to contiguous outages. The improvement, however, does not come from the physics term. The architecture-matched twin without physics is as accurate, and a well-tuned gradient-boosted tree on the same features is more accurate still in RMSE. What carries the gain is (i) the richer, mask-aware input representation, which lowers the baseline network's error by 14\% on its own, (ii) temporal context, which adds another 13\%, and (iii) mask awareness under outages, which is worth 10\% at 40\% block loss. The answer to the research question is therefore: \emph{physics-informed learning, as implemented here, improves the prediction only as part of a mask-aware temporal model; the physics term itself is neutral for accuracy and gives a modest reduction in heat-balance residual on complete data.}

\textbf{RQ1.} Random point loss up to 40\% has almost no effect on any model, because hourly building signals are smooth and the imputed value is close to the true one. \textbf{RQ2.} Contiguous multi-sensor outages are the real threat, raising RMSE by 15--23\% at 40\%; this confirms the argument of \citet{lee2024missing} that the missingness mechanism, not only the rate, decides the impact, and it means that robustness studies that only use MCAR masks overstate robustness. \textbf{RQ3.} The physics term lowers the RC residual by 10\% on complete data, but not under heavy outages, and it does not reduce implausible responses to outdoor temperature; the learned time constant is stable ($\tau\approx49$\,h) but not identifiable from the data.

\subsection{Why the physics prior did not help here}
The literature reports that physics priors help most when data are scarce or when models must extrapolate \citep{gokhale2022pinn,liguori2024opening}. Three features of this case weaken the prior. First, the indoor temperature is tightly controlled, so the temperature change carries little information about HVAC energy; a least-squares fit cannot even identify the envelope conductance. Second, the HVAC meter covers the rooftop units only, while fan-coil units and the central chilled- and hot-water plant also condition the space, so the balance of Eq.~\eqref{eq:rc} links the temperature to only part of the thermal input. Third, there are 2.5 years of hourly training data, and the network already learns the dominant relationships from data. This matches the open problems named by \citet{ma2025review}: the choice of physics prior and its match to what is metered decide whether PIML helps, and a mis-matched prior can only be down-weighted. The monotonic loss of accuracy as $\lambda$ grows is direct evidence of that mismatch. A physics-derived feature layer (fan affinity law and air-side sensible load) was also tried in a pilot with XGBoost and gave no consistent gain across seeds, so it was not adopted.

\subsection{Comparison with the literature}
The complete-data results agree with \citet{mahajan2024bnn} that their BNN is more accurate than an LSTM trained on the same per-hour inputs; they add that a recurrent network with a 24-hour window and mask features beats both. They agree with \citet{che2018grud} that exposing the missingness pattern to the model is valuable. They partly agree with \citet{liguori2024opening}: the physics term has little effect on accuracy with complete data, but unlike in their imputation study it does not buy graceful degradation for energy prediction. The strongest model, XGBoost, is consistent with the strong record of gradient-boosted trees on tabular prediction problems reported by \citet{chen2016xgboost}.

\subsection{Implications}
\textbf{For practitioners.} (1) Test energy models with contiguous multi-sensor outages, not just random gaps. (2) Give models the observation mask and the time since the last reading; it is cheap and it is the most effective single measure found here. (3) A gradient-boosted model with these features is a strong, fast default (under a second to train, 0.2\,µs per prediction); a recurrent model is preferable when typical-hour accuracy (MAE, MAPE) matters more than rare large errors. (4) Before adding a physics term, check that the metered energy and the modelled heat balance cover the same equipment.

\textbf{For researchers.} Physics-informed energy models should be reported with an architecture-matched $\lambda=0$ ablation; without it, the gain of the PI-GRU over the baseline would have been wrongly credited to physics. Identifiability of the physical parameters should be checked with a plain least-squares fit, because a stable learned parameter can be a product of the constraint rather than the data.

\subsection{Limitations}\label{sec:limits}
\begin{itemize}[nosep]
\item One building, one mild climate and one test period, which is also affected by COVID-19 occupancy changes; the conclusions about the physics prior may not transfer to buildings with a metered plant or larger temperature swings.
\item The published baseline numbers come from the paper; two of its inputs are not in the public release, so the re-implementation is close but not exact.
\item The RC model is first order and uses a mode proxy derived from rooftop-unit air temperatures; zone-level flows, chilled-water and hot-water energy are not available.
\item Masks are synthetic. They are designed to resemble real outages (contiguous, multi-channel), but real failure modes such as drift or stuck values are not simulated.
\item Five seeds give wide intervals for small effects; differences of about 1\% between the PI-GRU and its twin cannot be resolved.
\item Hyperparameter budgets are equal but small (five trials per family); larger searches could change the ranking of the recurrent models.
\end{itemize}
